\section*{الفصل \ref{ch:b3:sensory-systems} --- الأجهزة الحسية}

\begin{solution}{exo:b3:sensory-systems:1}
بالسلك لا بالرسالة: فكل عصب حسي \omterm{def:b3:sensory-systems:principles}{خط موسوم} ينتهي في منطقته هو
من الدماغ، وكل ما يصل على امتداد العصب البصري يُقرأ ضوءًا
--- ولذلك يُحدث الضغط على مقلة العين ومضةً وتُحدث ضربةٌ على
الأذن رنينًا.
\end{solution}

\begin{solution}{exo:b3:sensory-systems:2}
في الظلام يبقي GMP الحلقي القنوات الكاتيونية مفتوحة ويبقي
``تيار الظلام'' الداخل الثابت الخلية مزالة الاستقطاب ومطلِقة
للغلوتامات. ويفعّل الضوء \omterm{def:b3:sensory-systems:retina}{الرودوبسين} والترانسدوسين
والفسفودي إستراز، الذي يهدم cGMP؛ فتُغلق القنوات، ويتوقف
التيار الداخل، وتُفرط الخلية استقطابها وتطلق أقل --- فالضوء
يُشار إليه بنقصان.
\end{solution}

\begin{solution}{exo:b3:sensory-systems:3}
$I = I_{0}10^{L/10}$: فعند \qty{20}{dB}، \qty{1e-10}{W/m^2}؛ وعند
\qty{60}{dB}، \qty{1e-6}{W/m^2}؛ وعند \qty{110}{dB}، \qty{0.1}{W/m^2}.
ونسبة الأعلى إلى الأخفت: $10^{9}$.
\end{solution}

\begin{solution}{exo:b3:sensory-systems:4}
الحلو والأومامي والمر عبر مستقبلات مقترنة بالبروتين G
(واحد للحلو، وواحد للأومامي، ونحو خمسة وعشرين للمر)؛ والملح
عبر قناة الصوديوم ENaC؛ والحامض عبر قناة بروتونات. ومعظم
``النكهة'' هو رائحة جزيئات طيّارة تبلغ الأنف من الفم؛ وإذا
سُدّ الأنف لم تبقَ إلا الأذواق الخمسة وبدا الطعام بلا طعم.
\end{solution}

\begin{solution}{exo:b3:sensory-systems:5}
فِشنر: $S = (1/0.02)\ln(10\,000/100) = 50\ln 100 \approx 230$ خطوة.
وبالعدّ: $1.02^{n} = 100$ يعطي $n = \ln 100/\ln 1.02 = 233$.
\end{solution}

\begin{solution}{exo:b3:sensory-systems:6}
$P = 1 - e^{-a}\sum_{k=0}^{5}a^{k}/k!$: فعند $a = 2$: $0.017$؛ وعند $a = 6$:
$0.55$؛ وعند $a = 12$: $0.98$. وتُرى نصف الومضات عند $a \approx 5.7$؛
ومع امتصاص \qty{10}{\%} يجب أن توصل الومضة نحو $57$
فوتونًا عند القرنية.
\end{solution}

\begin{solution}{exo:b3:sensory-systems:7}
$55/3\times 1.3 \approx 24$؛ وبديسيبلات الضغط، $20\log_{10}24
\approx \qty{28}{dB}$. والعُظيمات الملتحمة لا تقدر على نقل
الاهتزاز، فلا يبلغ الصوت \omterm{def:b3:sensory-systems:cochlea}{القوقعة} إلا بالتوصيل عبر الجمجمة،
بلا كسب الأذن الوسطى: أي فقد توصيلي قدره نحو
\qtyrange{30}{50}{dB}.
\end{solution}

\begin{solution}{exo:b3:sensory-systems:8}
قناة \omterm{def:b3:sensory-systems:cochlea}{الخلية الشعرية} تفتحها \omterm{def:b3:sensory-systems:cochlea}{الوصلة القمية} بسحبها مباشرةً
--- بلا رسول، وبلا إنزيم --- فتستجيب في ميكروثوانٍ؛ وأما
المستقبِل الضوئي فيضخّم فوتونًا واحدًا عبر مرحلتين
إنزيميتين، وهما تستغرقان عشرات الميلي ثانية. فالسمع يكسب
دقة التوقيت اللازمة لتحديد موضع الأصوات بفروق الوصول بين
الأذنين ولتتبع الطور؛ والبصر يكسب الحساسية اللازمة لعدّ
الفوتونات المفردة، ويدفع من السرعة.
\end{solution}

\begin{solution}{exo:b3:sensory-systems:9}
$H(0.05) = -0.05\ln 0.05 - 0.95\ln 0.95 = 0.150 + 0.049 = 0.199$.
والإنسان: $\ln\binom{400}{20} \approx 400\times 0.199 = 80$، أي نحو
$10^{34}$ نمط. والفأر: $\ln\binom{1100}{55} \approx 1100\times 0.199
= 219$، أي نحو $10^{95}$.
\end{solution}

\begin{solution}{exo:b3:sensory-systems:10}
الوسطي $500\times 0.1/40 = 1.25$ حدث عفوي في النافذة. و$P(k
\le 5) = e^{-1.25}(1 + 1.25 + 0.78 + 0.33 + 0.10 + 0.03) = 0.998$، فيكون
$P(k \ge 6) \approx 2\times 10^{-3}$. ولو كانت العتبة واحدًا
لوقع حدث عفوي في معظم النوافذ ($1 - e^{-1.25} = 0.71$) ولكان
الظلام مملوءًا بالومضات؛ وعند ستة تقع الإنذارات الكاذبة مرة
في كل خمسمئة نافذة. فالعتبة موضوعة حيث يمليها ضجيج
العصيّات لا حساسيتها.
\end{solution}

\begin{solution}{exo:b3:sensory-systems:11}
أكبر تأخير $0.2/340 = \qty{590}{\micro s}$ (لصوت من أحد
الجانبين). وقرب الخط الأوسط $\Delta t \approx (d/c)\sin\theta$، فتقابل $\qty{10}{\micro
s}$ قيمة $\sin\theta = 0.017$، أي نحو \qty{1}{\degree}. وفوق
\qty{4}{kHz} يكون الدور (\qty{250}{\micro s}) أقصر مما تقدر
العصبونات على القفل عليه، ويمتد تأخير قدره بضع مئات من
الميكروثواني على أكثر من دورة، فيلتبس الطور؛ فيستعمل الدماغ
عندئذ فرق المستوى بين الأذنين وشفرة الموضع.
\end{solution}

\begin{solution}{exo:b3:sensory-systems:12}
بعيدًا عن الحد، $r = I(1 - 2w)$: أي $0.6$ في الجانب الداكن،
و$1.2$ في المضيء. وعند آخر مستقبل داكن ($I = 1$، وجاراه $1$
و$2$): $r
= 1 - 0.2\times 3 = 0.4$؛ وعند أول مستقبل مضيء: $r = 2 - 0.6 = 1.4$.
فتُبالَغ القفزة إلى هبوط وارتفاع --- وهي عصابات ماخ. وتكسب
\omterm{def:b3:sensory-systems:retina}{الشبكية} التباين والحواف، ومدًى ديناميًا أصغر تنقله؛ وتخسر
الوفاء للمستويات المطلقة، وتنتج أوهام سطوع.
\end{solution}

\begin{solution}{pb:b3:sensory-systems:1}
\textbf{1.} $0.01/3.6\times 10^{-19} = 2.8\times 10^{16}$ فوتون في
الثانية.
\textbf{2.} مساحة الحدقة $\pi(3.5\times 10^{-3})^{2} = 3.8\times 10^{-5}$
م$^{2}$. وعند \qty{1}{km}: النسبة $3.8\times 10^{-5}/1.26\times 10^{7} =
3\times 10^{-12}$، أي $8.5\times 10^{4}$ فوتون في الثانية. وعند \qty{10}{km}:
$850$ في الثانية.
\textbf{3.} في \qty{0.1}{s} مع امتصاص \qty{10}{\%}: $850$ عند
\qty{1}{km}، و$8.5$ عند \qty{10}{km} --- وكلاهما فوق ستة؛
فالشمعة تُرى على عشرة كيلومترات (في معظم الأحيان).
\textbf{4.} ستة ممتصّة في النافذة تحتاج إلى $600$ فوتون في
الثانية إلى الحدقة: أي $4\pi d^{2} = 3.8\times 10^{-5}\times 2.8\times 10^{16}/600
= 1.8\times 10^{9}$ م$^{2}$، و$d \approx \qty{12}{km}$. وهناك يكون $P(k \ge
6 \mid a = 6) = 0.55$: أي إن الشمعة تُرى في نحو نصف النوافذ.
\textbf{5.} $e^{-6} = 0.0025$. وعند العتبة يتذبذب العدد من
نافذة إلى نافذة --- $6 \pm 2.4$ --- فيبدو المنبع كأنه يأتي
ويذهب: وهو تلألؤ النجم الخافت (يضيف إليه الجو تلألؤه هو).
\textbf{6.} تذبذب الخلفية، $\sqrt{10^{4}} = 100$ فوتون في
النافذة، يغرق $8.5$ من الشمعة: فالكشف يقتضي أن تتجاوز
الإشارةُ ضجيجَ الخلفية، لا عتبةَ العين المظلمة.
\textbf{7.} $\qty{1}{pA}/200 = \qty{5}{fA}$ لكل قناة؛ و$5\times
10^{-15}/1.6\times 10^{-19} = 3\times 10^{4}$ أيون في الثانية.
\textbf{8.} جزء من ثلاثين؛ ونحو ثلاثين فوتونًا متزامنًا تغلق
كل قناة وتشبع \omterm{def:b3:sensory-systems:retina}{العصيّة}.
\textbf{9.} $10^{-12}\times 0.2 = 2\times 10^{-13}$ كولوم، أي $1.3\times 10^{6}$
كاتيون: ففوتون واحد يتحكم بأكثر من مليون شحنة.
\textbf{10.} تتشبع \omterm{def:b3:sensory-systems:retina}{العصيّة} في غضون ميلي ثوانٍ، وتُغلق كل
قناة، ويُقصَر رودوبسينها أسرع مما يُجدَّد؛ وأما المخاريط،
بكسبها الأقل وإطفائها الأسرع، فتبقى مستجيبة وتنقل مداها عبر
ست عشريات من الضوء.
\textbf{11.} الكسب العالي يحتاج إلى تكامل طويل (فكل مرحلة
تستغرق وقتًا وتُطال الاستجابة لتتراكم)، \omterm{def:b3:sensory-systems:retina}{فالعصيّة} بطيئة
حساسة؛ وأما تضخيم \omterm{def:b3:sensory-systems:retina}{المخروط} الأصغر وإطفاؤه الأسرع فيعطيان
استجابة في \qty{20}{ms}، سريعة بما يكفي للحركة، بثمن الحاجة
إلى مئة فوتون.
\textbf{12.} $10^{8}/40 = 2.5\times 10^{6}$ مثلنة عفوية في
الثانية عبر \omterm{def:b3:sensory-systems:retina}{الشبكية} --- غير أنها $1.25$ فقط في النافذة في
أي رقعة من $500$ \omterm{def:b3:sensory-systems:retina}{عصيّة}، أي أدنى بكثير من عتبة الستة، فيُرفض
الضجيج رقعةً رقعة؛ والعتبة والتجميع موجودان لهذا بالضبط.
\textbf{13.} \qty{0.1}{W/m^2}؛ وعلى \qty{5.5e-5}{m^2}، \qty{5.5e-6}{W}.
\textbf{14.} $5.5\times 10^{-6}\times 7200 = \qty{0.04}{J}$ --- أي نحو
أربعمئة قطرة مطر على مدى ساعتين.
\textbf{15.} $10^{-12}\times 5.5\times 10^{-5} = 5.5\times 10^{-17}$ واط؛
وفي \qty{10}{ms}، $5.5\times 10^{-19}$ جول --- أي نحو طاقة فوتون
مرئي واحد أو اثنين.
\textbf{16.} $0.3/5000 = 6\times 10^{-5}$ راديان، أي $0.003^{\circ}$؛
والانحراف نحو ثلاث ذرات هدروجين.
\textbf{17.} $10^{(110 - 85)/10} = 316$: فساعتان عند \qty{110}{dB}
توصلان طاقة $630$ ساعة عند \qty{85}{dB} --- أي قرابة شهر من
أيام العمل.
\textbf{18.} نحو $120$ خطوة ملحوظة بالكاد، واحدة لكل \omterm{def:b3:sensory-systems:cochlea}{ديسيبل}.
\textbf{19.} $2^{400} = 10^{400\log_{10}2} = 10^{120}$.
\textbf{20.} $\ln\binom{400}{20} \approx 400\,H(0.05) \approx 80$، ناقصًا
تصحيح ستيرلنغ قدره نحو $2$: أي نحو $10^{34}$ نمط.
\textbf{21.} عُدّ المستقبِلات التي في أحد النمطين دون الآخر:
$5 +
5 = 10$ من $40$ (وهي مسافة هامنغ)؛ والأنماط المتداخلة بثلاثة
أرباعها تسوق الكبيبات نفسها تقريبًا وتُقرأ رائحةً واحدة
تقريبًا.
\textbf{22.} $10^{12}/10^{34} = 10^{-22}$ من الأنماط. فليست
الشفرة هي الحد: بل المستقبِلات مشوَّشة، وكثير منها منغَّم
متشابهًا فترتبط فعالياتها، والجزيئات الحقيقية لا تنتج إلا
مجموعة جزئية صغيرة من الأنماط، وتمييز الدماغ للأنماط
وذاكرته لها محدودان.
\textbf{23.} $\ln\binom{1100}{20} \approx 1100\,H(0.018) = 1100\times
0.091 \approx 100$، أي نحو $10^{43}$ --- أي نحو $10^{9}$ ضعف عدد
أنماط العشرين مستقبلًا عند الإنسان (وأكثر بكثير إن استعمل
الفأر مستقبِلات أكثر لكل رائحة).
\textbf{24.} الشم عمومًا لا يكاد يتغير --- فالشفرة مكرَّرة
ومستقبِلات أخرى تحمل كل جزيء رائحة --- غير أن جزيء رائحة
يعتمد على ذلك المستقبل عند تركيز منخفض يصير غير قابل للكشف
أو يتبدل طابعه: أي فقد شم نوعي، كما عند الكثيرين الذين لا
يشمّون الأندروستينون.
\textbf{25.} تُرى الشمعة إلى نحو \qty{12}{km}؛ ويوصل صوت
العتبة $5.5\times 10^{-19}$ جول إلى طبلة الأذن في \qty{10}{ms}،
أي فوتونًا أو فوتونين؛ ويقدر أنف بشري على تكوين نحو
$10^{34}$ نمط من عشرين مستقبلًا.
\end{solution}
